\section*{Chapitre \ref{ch:g11:igneous-rocks} --- Les roches magmatiques}

\begin{solution}{exo:g11:igneous-rocks:1}
Une \omterm{def:g11:igneous-rocks:volcanic-plutonic}{roche volcanique} s'est formée à partir d'un magma solidifié en surface,
sous forme de lave~; une \omterm{def:g11:igneous-rocks:volcanic-plutonic}{roche plutonique} a cristallisé en profondeur, à
l'intérieur d'autres roches. Les paires~: basalte et \omterm{def:g11:ocean-ridges:gabbro}{gabbro} (\omterm{def:g11:igneous-rocks:felsic-mafic}{mafiques}),
\omterm{def:g11:igneous-rocks:rock-names}{andésite} et \omterm{def:g11:igneous-rocks:rock-names}{diorite} (\omterm{def:g11:igneous-rocks:felsic-mafic}{intermédiaires}), \omterm{def:g11:igneous-rocks:rock-names}{rhyolite} et granite (\omterm{def:g11:igneous-rocks:felsic-mafic}{felsiques}).
\end{solution}

\begin{solution}{exo:g11:igneous-rocks:2}
(1) Vitreuse. (2) À gros grains. (3) \omterm{def:g11:igneous-rocks:porphyritic}{Porphyrique} (\omterm{def:g11:igneous-rocks:porphyritic}{phénocristaux} d'olivine
dans une pâte fine). (4) À grain fin.
\end{solution}

\begin{solution}{exo:g11:igneous-rocks:3}
\qty{48}{\percent}~: \omterm{def:g11:igneous-rocks:felsic-mafic}{mafique}, \omterm{def:g11:ocean-ridges:gabbro}{gabbro}. \qty{60}{\percent}~: \omterm{def:g11:igneous-rocks:felsic-mafic}{intermédiaire},
\omterm{def:g11:igneous-rocks:rock-names}{diorite}. \qty{72}{\percent}~: \omterm{def:g11:igneous-rocks:felsic-mafic}{felsique}, granite.
\end{solution}

\begin{solution}{exo:g11:igneous-rocks:4}
Un \omterm{def:g11:igneous-rocks:intrusion}{dyke} est une lame de roche magmatique qui recoupe les couches des roches
qu'il traverse~; un \omterm{def:g11:igneous-rocks:intrusion}{sill} est une lame parallèle aux couches, insérée entre
elles~; un \omterm{def:g11:igneous-rocks:intrusion}{pluton} est une grande \omterm{def:g11:igneous-rocks:intrusion}{intrusion} de forme quelconque, formée en
profondeur. C'est le \omterm{def:g11:igneous-rocks:intrusion}{dyke} vertical issu du \omterm{def:g11:igneous-rocks:intrusion}{pluton} qui a alimenté le volcan.
\end{solution}

\begin{solution}{exo:g11:igneous-rocks:5}
L'olivine se forme en premier~; le feldspath potassique, le quartz et
parfois la muscovite en dernier. Le plagioclase est d'abord riche en
calcium, puis de plus en plus riche en sodium à mesure que le magma
refroidit.
\end{solution}

\begin{solution}{exo:g11:igneous-rocks:6}
(1) A~: \qty{74}{\percent} de \ce{SiO2}, \omterm{def:g11:igneous-rocks:felsic-mafic}{felsique}~; B~: \qty{47}{\percent},
\omterm{def:g11:igneous-rocks:felsic-mafic}{mafique}. (2) A est pauvre en \ce{FeO}, \ce{MgO} et \ce{CaO} et riche en
\ce{Na2O} et \ce{K2O}, comme un magma \omterm{def:g11:igneous-rocks:felsic-mafic}{felsique}~; B est riche en \ce{FeO},
\ce{MgO} et \ce{CaO} et pauvre en \ce{Na2O}, comme un magma \omterm{def:g11:igneous-rocks:felsic-mafic}{mafique}. (3) A
est une \omterm{def:g11:igneous-rocks:rock-names}{rhyolite}~; B un \omterm{def:g11:ocean-ridges:gabbro}{gabbro}.
\end{solution}

\begin{solution}{exo:g11:igneous-rocks:7}
La durée croît comme le carré de l'épaisseur. \qty{1}{m}~: $10/4 = 2.5$
jours. \qty{20}{m}~: $10 \times 10^2 = 1000$ jours, environ 2.7 ans.
\qty{200}{m}~: $10 \times 100^2 = \num{100000}$ jours, environ 270 ans.
\end{solution}

\begin{solution}{exo:g11:igneous-rocks:8}
(1) Des \unit{m^2} divisés par des \unit{m^2/s} donnent des secondes.
(2) $400/10^{-6} =
\qty{4e8}{s}$, soit $4 \times 10^8 / 3.16 \times 10^7 \approx 12.7$ ans. (3) $0.21 \times 12.7 \approx 2.7$ ans.
\end{solution}

\begin{solution}{exo:g11:igneous-rocks:9}
Elle est \omterm{def:g11:igneous-rocks:porphyritic}{porphyrique}~: des \omterm{def:g11:igneous-rocks:porphyritic}{phénocristaux} de plagioclase dans une pâte à
grain fin. Le magma a d'abord refroidi lentement en profondeur, et les
cristaux de plagioclase sont devenus grands tandis que l'essentiel du magma
restait liquide~; puis il est monté et a refroidi vite, et le liquide
restant a cristallisé en une pâte fine.
\end{solution}

\begin{solution}{exo:g11:igneous-rocks:10}
(1) À gauche~: 2 Mg, 2 Si, 6 O~; à droite~: 2 Mg, 2 Si, 6 O. (2) \qty{1}{kg}
d'olivine représente $1000/140.7 \approx 7.11$~mol, qui prennent autant de
moles de silice~: $7.11 \times 60.1 \approx \qty{427}{g}$. (3) Quand il ne
reste plus de silice dans le liquide, la réaction ne peut pas continuer~:
l'olivine reste.
\end{solution}

\begin{solution}{exo:g11:igneous-rocks:11}
(1) $2.5 \times 20 \times 1100 = \qty{55000}{MJ}$, soit \qty{5.5e10}{J}.
(2) $900 \times 20 = \qty{18000}{MJ}$, soit \qty{1.8e10}{J}. (3) La chaleur
de cristallisation vaut environ le tiers de la chaleur perdue par
refroidissement, le quart du total~: assez pour ralentir beaucoup la
solidification, car elle est entièrement libérée au-dessus du solidus.
\end{solution}

\begin{solution}{exo:g11:igneous-rocks:12}
(1) Les plus petits cristaux au contact, qui s'est figé en moins d'un
jour~; les plus grands au centre, resté liquide pendant environ 950 jours.
(2) Non~: les \qty{5}{m} extérieurs ont demandé environ 545 jours, les
\qty{5}{m} suivants, jusqu'au centre, environ 400 jours~; le front
ralentit d'abord, puis accélère près du centre, là où se rejoignent les
pertes de chaleur par les deux côtés. (3) Même en son centre, le \omterm{def:g11:igneous-rocks:intrusion}{sill}
s'est figé en moins de trois ans~; un granite aux grands cristaux demande
un temps bien plus long.
\end{solution}

\begin{solution}{exo:g11:igneous-rocks:13}
(1) Le \omterm{def:g11:igneous-rocks:intrusion}{pluton} est 100 fois plus épais~: $2.6 \times 100^2 = \num{26000}$
ans. (2) Par exemple~: un \omterm{def:g11:igneous-rocks:intrusion}{pluton} profond se trouve dans une roche déjà
chaude, et perd donc sa chaleur plus lentement que dans une roche à
\qty{50}{\degreeCelsius}~; il peut être alimenté par du magma neuf encore
et encore~; un \omterm{def:g11:igneous-rocks:intrusion}{pluton} n'est pas une lame infinie~; les propriétés des
roches changent avec la température.
\end{solution}

\begin{solution}{exo:g11:igneous-rocks:14}
Un magma \omterm{def:g11:igneous-rocks:felsic-mafic}{mafique} est chaud et pauvre en silice~: il cristallise les
premiers minéraux de la série --- olivine, pyroxène, plagioclase riche en
calcium --- et manque de silice avant que le quartz puisse se former. Un
magma \omterm{def:g11:igneous-rocks:felsic-mafic}{felsique} est moins chaud et riche en silice~: l'olivine n'y est pas
stable, et il va jusqu'au bout de la série, en formant un plagioclase riche
en sodium, du feldspath potassique et du quartz.
\end{solution}

\begin{solution}{exo:g11:igneous-rocks:15}
(1) $b = \sqrt{t\kappa/C} = \sqrt{86400 \times 10^{-6}/0.21} \approx
\qty{0.64}{m}$.
(2) La loi donne $0.21 \times 25/10^{-6} \approx
\qty{5.3e6}{s}$, environ
deux mois~: le même ordre de grandeur qu'un mois (une coulée de lave perd
aussi de la chaleur dans l'air, par le haut). (3) Vitreuse pour les
matériaux projetés dans l'air ou dans l'eau~; à grain fin pour les coulées
de lave et les \omterm{def:g11:minerals:thin-section}{lames minces}~; à grain plus gros, mais encore petit, pour
les \omterm{def:g11:igneous-rocks:intrusion}{intrusions} peu profondes épaisses.
\end{solution}

\begin{solution}{pb:g11:igneous-rocks:1}
\textbf{1.} Un \omterm{def:g11:igneous-rocks:intrusion}{sill}~: il est parallèle aux bancs, inséré entre eux.

\textbf{2.} Au contact, à grain très fin, presque vitreuse~; au centre, à
grain fin à moyen, les minéraux tout juste visibles à la loupe.

\textbf{3.} Le contact a perdu sa chaleur le plus vite, dans le grès
froid~: les cristaux n'ont pas eu le temps de grandir.

\textbf{4.} Une bordure figée. Elle prouve que le magma était chaud quand
il s'est mis en place, et le grès froid.

\textbf{5.} Un \omterm{def:g11:igneous-rocks:intrusion}{dyke}, le filon d'alimentation du \omterm{def:g11:igneous-rocks:intrusion}{sill}.

\textbf{6.} \omterm{def:g11:igneous-rocks:felsic-mafic}{Mafique}~: \qty{49}{\percent} de \ce{SiO2}, riche en fer, en
magnésium et en calcium.

\textbf{7.} De l'olivine, du pyroxène et un plagioclase riche en calcium.

\textbf{8.} Un basalte~; la \omterm{def:g11:igneous-rocks:volcanic-plutonic}{roche plutonique} de même composition est le
\omterm{def:g11:ocean-ridges:gabbro}{gabbro}.

\textbf{9.} La \omterm{def:g11:igneous-rocks:texture}{texture} dépend de la vitesse de refroidissement, qui varie du
contact au centre, et non de la composition.

\textbf{10.} L'olivine d'abord~; le plagioclase commence à peu près en même
temps que le pyroxène.

\textbf{11.} $b^2/\kappa = 400/10^{-6} = \qty{4e8}{s}$.

\textbf{12.} $0.206 \times 4 \times 10^8 \approx \qty{8.2e7}{s}$, soit
$8.2 \times 10^7 / 3.16 \times 10^7 \approx 2.6$ ans.

\textbf{13.} $0.206 \times 4/10^{-6} \approx \qty{8.2e5}{s}$, environ
9.5~jours.

\textbf{14.} $0.07 \times 4 \times 10^8 = \qty{2.8e7}{s}$, environ 0.9 an.
Le magma qui cristallise libère de la chaleur, qu'il faut évacuer elle
aussi avant qu'il devienne solide.

\textbf{15.} Environ 950 jours contre 0.6 jour~: environ 1600 fois plus
vite.

\textbf{16.} Environ 950 jours, soit quelque 2.6 ans~: bien plus qu'une coulée
de lave, dans les mois ou les années des \omterm{def:g11:igneous-rocks:intrusion}{intrusions} peu profondes.

\textbf{17.} Mille fois plus de temps demande une épaisseur $\sqrt{1000} \approx
32$ fois plus grande~: environ \qty{630}{m}.

\textbf{18.} Près de la surface, un corps de magma ferait éruption, ou
refroidirait trop vite~; un grand corps qui refroidit pendant très
longtemps doit être profond, dans une roche déjà chaude.

\textbf{19.} Elles ont été usées par l'érosion.

\textbf{20.} Le \omterm{def:g11:igneous-rocks:intrusion}{sill} a mis environ \textbf{2.6 ans} à se solidifier.
\end{solution}
